%% ma: L12-B08-0172
%% lop: 12
%% chuong: 2
%% bai: 8
%% dang: 12.8.53
%% loai: TN4
%% muc: VD
%% dap_an: "A"
%% nguon: "CĐ2.4.pdf giáo viên gửi 27/10/2026 (M07, Lê Minh Khôi – Nguyễn Anh Khôi, Chương II Chủ đề 4 Các bài toán về tâm tỉ cự), Câu 22"
%% kien_thuc: "Bài 8 (khoảng cách hai điểm, biểu thức tọa độ), Bài 6 (bình phương vô hướng, tâm tỉ cự)"
%% trang_thai: chinh-thuc
%% ghi_chu: "Giáo viên đã duyệt 27/10/2026. Đề gốc dùng chung giả thiết cho Câu 20–22; đã chép vào từng câu"
\begin{ex}
Trong không gian $Oxyz$, cho ba điểm $M(2;1;-2)$, $N(-1;0;3)$, $P(1;-3;5)$ và điểm $I$ bất kỳ trong không gian. Biết rằng $I$ thuộc trục tung, tìm giá trị lớn nhất của biểu thức $V=IM^2-5IN^2+IP^2$.
\choice{\True $V_{\max}=-\dfrac{14}3$}{$V_{\max}=-\dfrac{16}3$}{$V_{\max}=\dfrac{16}3$}{$V_{\max}=\dfrac{14}3$}
\loigiai{$I(0;y;0)$: $V=(4+(y-1)^2+4)-5(1+y^2+9)+(1+(y+3)^2+25)=-3y^2+4y-6$.

$V$ lớn nhất khi $y=\dfrac23$: $V_{\max}=-\dfrac43+\dfrac83-6=-\dfrac{14}3$.}
\end{ex}
